\documentclass[a4paper,10pt]{paper}
%DIF LATEXDIFF DIFFERENCE FILE
%DIF DEL solution2.tex           Mon Sep 29 04:32:56 2025
%DIF ADD revised/solution2.tex   Thu Sep  3 07:04:33 2026

\usepackage{mycommands}
%DIF 4d4
%DIF < \usepackage{amsmath}
%DIF -------

\title{Seismology Problem Set 2 \DIFdelbegin \DIFdel{- }\DIFdelend \DIFaddbegin \DIFadd{-- }\DIFaddend Solutions}
\author{David Al-Attar \\ Michaelmas Term, \the\year}
%DIF PREAMBLE EXTENSION ADDED BY LATEXDIFF
%DIF UNDERLINE PREAMBLE %DIF PREAMBLE
\RequirePackage[normalem]{ulem} %DIF PREAMBLE
\RequirePackage{color}\definecolor{RED}{rgb}{1,0,0}\definecolor{BLUE}{rgb}{0,0,1} %DIF PREAMBLE
\providecommand{\DIFadd}[1]{{\protect\color{blue}\uwave{#1}}} %DIF PREAMBLE
\providecommand{\DIFdel}[1]{{\protect\color{red}\sout{#1}}}                      %DIF PREAMBLE
%DIF SAFE PREAMBLE %DIF PREAMBLE
\providecommand{\DIFaddbegin}{} %DIF PREAMBLE
\providecommand{\DIFaddend}{} %DIF PREAMBLE
\providecommand{\DIFdelbegin}{} %DIF PREAMBLE
\providecommand{\DIFdelend}{} %DIF PREAMBLE
\providecommand{\DIFmodbegin}{} %DIF PREAMBLE
\providecommand{\DIFmodend}{} %DIF PREAMBLE
%DIF FLOATSAFE PREAMBLE %DIF PREAMBLE
\providecommand{\DIFaddFL}[1]{\DIFadd{#1}} %DIF PREAMBLE
\providecommand{\DIFdelFL}[1]{\DIFdel{#1}} %DIF PREAMBLE
\providecommand{\DIFaddbeginFL}{} %DIF PREAMBLE
\providecommand{\DIFaddendFL}{} %DIF PREAMBLE
\providecommand{\DIFdelbeginFL}{} %DIF PREAMBLE
\providecommand{\DIFdelendFL}{} %DIF PREAMBLE
%DIF COLORLISTINGS PREAMBLE %DIF PREAMBLE
\RequirePackage{listings} %DIF PREAMBLE
\RequirePackage{color} %DIF PREAMBLE
\lstdefinelanguage{DIFcode}{ %DIF PREAMBLE
%DIF DIFCODE_UNDERLINE %DIF PREAMBLE
  moredelim=[il][\color{red}\sout]{\%DIF\ <\ }, %DIF PREAMBLE
  moredelim=[il][\color{blue}\uwave]{\%DIF\ >\ } %DIF PREAMBLE
} %DIF PREAMBLE
\lstdefinestyle{DIFverbatimstyle}{ %DIF PREAMBLE
	language=DIFcode, %DIF PREAMBLE
	basicstyle=\ttfamily, %DIF PREAMBLE
	columns=fullflexible, %DIF PREAMBLE
	keepspaces=true %DIF PREAMBLE
} %DIF PREAMBLE
\lstnewenvironment{DIFverbatim}{\lstset{style=DIFverbatimstyle}}{} %DIF PREAMBLE
\lstnewenvironment{DIFverbatim*}{\lstset{style=DIFverbatimstyle,showspaces=true}}{} %DIF PREAMBLE
%DIF END PREAMBLE EXTENSION ADDED BY LATEXDIFF

\begin{document}

\maketitle

\begin{enumerate}
\item From Lecture\DIFaddbegin \DIFadd{~}\DIFaddend 18, we know that the least squares solution for the given problem is
\DIFdelbegin \begin{displaymath}
    \DIFdel{\mathbf{m}=(\mathbf{A}^{T}\mathbf{C}^{-1}\mathbf{A}+\lambda \mathbf{B})^{-1}\mathbf{A}^{T}\mathbf{C}^{-1}\mathbf{d}
}\end{displaymath}%DIFAUXCMD
\DIFdelend \DIFaddbegin \begin{equation*}
    \DIFadd{\mathbf{m}=(\mathbf{A}^{T}\mathbf{C}^{-1}\mathbf{A}+\lambda \mathbf{B})^{-1}\mathbf{A}^{T}\mathbf{C}^{-1}\mathbf{d},
}\end{equation*}\DIFaddend 
with $\mathbf{d}$ the data vector. Here we take \DIFdelbegin \DIFdel{$\mathbf{d}=\mathbf{A}\mathbf{m}_{in}$}\DIFdelend \DIFaddbegin \DIFadd{$\mathbf{d}=\mathbf{A}\mathbf{m}_{\mathrm{in}}$}\DIFaddend , and hence find that
\DIFdelbegin \begin{displaymath}
    \DIFdel{\mathbf{m}_{out}=(\mathbf{A}^{T}\mathbf{C}^{-1}\mathbf{A}+\lambda \mathbf{B})^{-1}\mathbf{A}^{T}\mathbf{C}^{-1}\mathbf{A}\mathbf{m}_{in}.
}\end{displaymath}%DIFAUXCMD
\DIFdelend \DIFaddbegin \begin{equation*}
    \DIFadd{\mathbf{m}_{\mathrm{out}}=(\mathbf{A}^{T}\mathbf{C}^{-1}\mathbf{A}+\lambda \mathbf{B})^{-1}\mathbf{A}^{T}\mathbf{C}^{-1}\mathbf{A}\mathbf{m}_{\mathrm{in}}.
}\end{equation*}\DIFaddend 
It follows that the resolution matrix takes the form
\begin{equation*}
    \mathbf{R}=(\mathbf{A}^{T}\mathbf{C}^{-1}\mathbf{A}+\lambda \mathbf{B})^{-1}\mathbf{A}^{T}\mathbf{C}^{-1}\mathbf{A}.
\end{equation*}
The resolution matrix relates the input and output of a synthetic inversion that does not include data errors. In an ideal case, we would like the resulting model to be equal to the input, and hence we want $\mathbf{R}$ to be close to the identity matrix.

\item
\begin{enumerate}
    \item The condition is that $\|\mathbf{A}\mathbf{m}-\mathbf{d}\|^{2} \le \chi^{2}_{c}$. Expanding this, we have
    \begin{equation*}
        (\mathbf{A}\mathbf{m}-\mathbf{d})^{T}(\mathbf{A}\mathbf{m}-\mathbf{d}) = \mathbf{m}^{T}\mathbf{A}^{T}\mathbf{A}\mathbf{m}-2\mathbf{m}^{T}\mathbf{A}^{T}\mathbf{d}+\mathbf{d}^{T}\mathbf{d} \le \chi_{c}^{2}.
    \end{equation*}
    Since the number of model parameters is greater than the number of data points ($m>n$), the problem is underdetermined. This means the $n \times m$ matrix $\mathbf{A}$ necessarily has a non-trivial null space.
     Consequently, the $m \times m$ matrix $\mathbf{G} = \mathbf{A}^{T}\mathbf{A}$ is positive semi-definite but \emph{singular} (not invertible). The singularity implies that the misfit is insensitive to any component
     of $\mathbf{m}$ that lies in the null space of $\mathbf{A}$. The region defined by the inequality is therefore unbounded in these \DIFdelbegin \DIFdel{null space }\DIFdelend \DIFaddbegin \DIFadd{null-space }\DIFaddend directions. This type of unbounded, ellipse-like shape is a
     \emph{hyper-cylinder} (or more precisely, an ellipsoidal cylinder).

    \item We want to minimise $\|\mathbf{m}\|^{2}$ subject to the constraint $\|\mathbf{A}\mathbf{m}-\mathbf{d}\|^{2} = \chi^{2}_{c}$. The Lagrangian is
    \begin{equation*}
        \mathcal{L}(\mathbf{m},\mu)=\mathbf{m}^{T}\mathbf{m}+\mu(\|\mathbf{A}\mathbf{m}-\mathbf{d}\|^{2}-\chi^{2}_{c}).
    \end{equation*}
    To find the minimum, we set the derivative with respect to $\mathbf{m}$ to zero:
    \begin{equation*}
        \frac{\partial\mathcal{L}}{\partial\mathbf{m}} = 2\mathbf{m}+2\mu\mathbf{A}^{T}(\mathbf{A}\mathbf{m}-\mathbf{d}) = \mathbf{0}.
    \end{equation*}
    Rearranging this gives
    \DIFdelbegin \begin{displaymath}
        \DIFdel{(\mathbf{I}+\mu\mathbf{A}^{T}\mathbf{A})\mathbf{m}=2\mu\mathbf{A}^{T}\mathbf{d},
    }\end{displaymath}%DIFAUXCMD
\DIFdelend \DIFaddbegin \begin{equation*}
        \DIFadd{(\mathbf{I}+\mu\mathbf{A}^{T}\mathbf{A})\mathbf{m}=\mu\mathbf{A}^{T}\mathbf{d},
    }\end{equation*}\DIFaddend 
    which can be solved for $\mathbf{m}$ to yield
    \DIFdelbegin \begin{displaymath}
        \DIFdel{\mathbf{m}=( \mathbf{A}^{T}\mathbf{A} + \frac{1}{\mu}\mathbf{I})^{-1}\mathbf{A}^{T}\mathbf{d}.
    }\end{displaymath}%DIFAUXCMD
\DIFdelend \DIFaddbegin \begin{equation*}
        \DIFadd{\mathbf{m}=\left( \mathbf{A}^{T}\mathbf{A} + \frac{1}{\mu}\mathbf{I}\right)^{-1}\mathbf{A}^{T}\mathbf{d}.
    }\end{equation*}\DIFaddend 
    The Lagrange multiplier $\mu$ is then found by substituting this expression for $\mathbf{m}$ back into the constraint equation, $\|\mathbf{A}\mathbf{m}-\mathbf{d}\|^{2} = \chi^{2}_{c}$.

    This result is identical in form to the \emph{regularised least squares solution}
    \DIFdelbegin \begin{displaymath}
        \DIFdel{\mathbf{m}_{reg}=(\mathbf{A}^{T}\mathbf{A}+\lambda \mathbf{I})^{-1}\mathbf{A}^{T}\mathbf{d},
    }\end{displaymath}%DIFAUXCMD
\DIFdelend \DIFaddbegin \begin{equation*}
        \DIFadd{\mathbf{m}_{\mathrm{reg}}=(\mathbf{A}^{T}\mathbf{A}+\lambda \mathbf{I})^{-1}\mathbf{A}^{T}\mathbf{d},
    }\end{equation*}\DIFaddend 
    if we identify the regularisation parameter $\lambda$ with $1/\mu$. The difference lies in how the regularisation parameter is chosen. In standard regularisation, $\lambda$ is often chosen subjectively. In this approach, however, the
     parameter is determined objectively by the requirement that the misfit must equal a statistically determined value, $\chi_{c}^{2}$. This is not, however, to say that the resulting model is the correct one. But it does have
     the potentially desirable property that of all the models that fit the data it has the smallest norm.
\end{enumerate}




\item A suitable Lagrangian for this problem can be defined by
\begin{equation*}
    L=J(\mathbf{u})-\int_{M}\frac{\partial}{\partial x_{j}}\left(A_{ijkl}\frac{\partial u_{k}}{\partial x_{l}}\right)u_{i}'\dd^{3}\mathbf{x}+\int_{\partial M}\left[\hat{n}_{j}\left(A_{ijkl}\frac{\partial u_{k}}{\partial x_{l}}\right)-\sigma g_{i}\right]w_{i}'\dd S,
\end{equation*}
where we have used the elastic tensor $A_{ijkl}=\lambda\delta_{ij}\delta_{kl}+\mu(\delta_{ik}\delta_{jl}+\delta_{il}\delta_{jk})$, and $u_{i}'$ and $w_{i}'$ are Lagrange multipliers. As usual, variation \DIFdelbegin \DIFdel{w.r.t }\DIFdelend \DIFaddbegin \DIFadd{with respect }\DIFaddend to the multipliers just gives the forward problem. Varying $u_{i}$, however, leads to the new condition
\begin{equation*}
    \int_{\partial M}h_{i}'\delta u_{i}\dd S-\int_{M}\frac{\partial}{\partial x_{j}}\left(A_{ijkl}\frac{\partial\delta u_{k}}{\partial x_{l}}\right)u_{i}'\dd^{3}\mathbf{x}+\int_{\partial M}\hat{n}_{j}\left(A_{ijkl}\frac{\partial\delta u_{k}}{\partial x_{l}}\right)w_{i}'\dd S=0,
\end{equation*}
where we have defined
\DIFdelbegin \begin{displaymath}
    \DIFdel{\mathbf{h}'(\mathbf{x})=\sum_{i=1}^{m}\frac{1}{\sigma_{i}^{2}}[\mathbf{u}(\mathbf{x}_{i})-\mathbf{u}_{i}^{obs}]\delta(\mathbf{x}-\mathbf{x}_{i}),
}\end{displaymath}%DIFAUXCMD
\DIFdelend \DIFaddbegin \begin{equation*}
    \DIFadd{\mathbf{h}'(\mathbf{x})=\sum_{i=1}^{m}\frac{1}{\varepsilon_{i}^{2}}[\mathbf{u}(\mathbf{x}_{i})-\mathbf{u}_{i}^{\mathrm{obs}}]\delta(\mathbf{x}-\mathbf{x}_{i}),
}\end{equation*}\DIFaddend 
to simplify the contribution associated with $J(\mathbf{u})$. To proceed, we need to integrate by parts to isolate $\delta u_{i}$. The key step is
\begin{align*}
    -\int_{M}&\frac{\partial}{\partial x_{j}}\left(A_{ijkl}\frac{\partial\delta u_{k}}{\partial x_{l}}\right)u_{i}'\dd^{3}\mathbf{x} \\
    &=-\int_{M}\frac{\partial}{\partial x_{j}}\left(A_{ijkl}\frac{\partial\delta u_{k}}{\partial x_{l}}u_{i}'\right)\dd^{3}\mathbf{x}+\int_{M}A_{ijkl}\frac{\partial\delta u_{k}}{\partial x_{l}}\frac{\partial u_{i}'}{\partial x_{j}}\dd^{3}\mathbf{x} \\
    &=-\int_{\partial M}\hat{n}_{j}\left(A_{ijkl}\frac{\partial\delta u_{k}}{\partial x_{l}}\right)u_{i}'\dd S+\int_{M}A_{klij}\frac{\partial\delta u_{i}}{\partial x_{j}}\frac{\partial u_{k}'}{\partial x_{l}}\dd^{3}\mathbf{x} \\
    &=-\int_{\partial M}\hat{n}_{j}\left(A_{ijkl}\frac{\partial\delta u_{k}}{\partial x_{l}}\right)u_{i}'\dd S+\int_{M}\frac{\partial}{\partial x_{j}}\left(A_{ijkl}\delta u_{i}\frac{\partial u_{k}'}{\partial x_{l}}\right)\dd^{3}\mathbf{x} -\int_{M}\delta u_{i}\frac{\partial}{\partial x_{j}}\left(A_{ijkl}\frac{\partial u_{k}'}{\partial x_{l}}\right)\dd^{3}\mathbf{x} \\
    &=-\int_{\partial M}\hat{n}_{j}\left(A_{ijkl}\frac{\partial\delta u_{k}}{\partial x_{l}}\right)u_{i}'\dd S+\int_{\partial M}\delta u_{i}\hat{n}_{j}\left(A_{ijkl}\frac{\partial u_{k}'}{\partial x_{l}}\right)\dd S -\int_{M}\delta u_{i}\frac{\partial}{\partial x_{j}}\left(A_{ijkl}\frac{\partial u_{k}'}{\partial x_{l}}\right)\dd^{3}\mathbf{x},
\end{align*}
where we note that the hyperelastic symmetry $A_{ijkl}=A_{klij}$ has been used. Using this result, the condition now reads
\begin{equation*}
    \int_{\partial M}[h_{i}'+\hat{n}_{j}\left(A_{ijkl}\frac{\partial u_{k}'}{\partial x_{l}}\right)]\delta u_{i}\dd S-\int_{M}\delta u_{i}\frac{\partial}{\partial x_{j}}\left(A_{ijkl}\frac{\partial u_{k}'}{\partial x_{l}}\right)\dd^{3}\mathbf{x} +\int_{\partial M}\hat{n}_{j}\left(A_{ijkl}\frac{\partial\delta u_{k}}{\partial x_{l}}\right)(w_{i}'-u_{i}')\dd S=0.
\end{equation*}
It follows that $u_{i}'$ must satisfy
\begin{equation*}
    \frac{\partial}{\partial x_{j}}\left(A_{ijkl}\frac{\partial u_{k}'}{\partial x_{l}}\right)=0,
\end{equation*}
subject to the boundary condition
\begin{equation*}
    \hat{n}_{j}\left(A_{ijkl}\frac{\partial u_{k}'}{\partial x_{l}}\right)=-h_{i}',
\end{equation*}
while also $w_{i}'=u_{i}'$ on $\partial M$. These are the adjoint equations for this problem. Having solved the forward and adjoint problems, the first-order perturbation to $J$ can be written
\begin{equation*}
    \delta J = -\int_{M}\frac{\partial}{\partial x_{j}}\left(\delta A_{ijkl}\frac{\partial u_{k}}{\partial x_{l}}\right)u_{i}'\dd^{3}\mathbf{x} + \int_{\partial M}\left[\hat{n}_{j}\left(\delta A_{ijkl}\frac{\partial u_{k}}{\partial x_{l}}\right)-\delta\sigma g_{i}\right]u_{i}'\dd S.
\end{equation*}
Performing an integration by parts, this simplifies to
\begin{equation*}
    \delta J=\int_{M}\delta A_{ijkl}\frac{\partial u_{k}}{\partial x_{l}}\frac{\partial u_{i}'}{\partial x_{j}}\dd^{3}\mathbf{x}-\int_{\partial M}\delta\sigma g_{i}u_{i}'\dd S.
\end{equation*}
From this we see that the sensitivity kernel \DIFdelbegin \DIFdel{w.r.t }\DIFdelend \DIFaddbegin \DIFadd{with respect to }\DIFaddend $\sigma$ is equal to $K_{\sigma}=-g_{i}u_{i}'$. To get the kernels for $\lambda$ and $\mu$, we note that
\begin{equation*}
    \delta A_{ijkl}=\delta\lambda\delta_{ij}\delta_{kl}+\delta\mu(\delta_{ik}\delta_{jl}+\delta_{il}\delta_{jk}),
\end{equation*}
from which we find
\DIFdelbegin \begin{displaymath}
    \DIFdel{\delta J=\int_{M}\delta\lambda\frac{\partial u_{i}}{\partial x_{i}}\frac{\partial u_{j}'}{\partial x_{j}}\dd^{3}\mathbf{x}+\int_{M}\delta\mu\left(\frac{\partial u_{i}}{\partial x_{j}}+\frac{\partial u_{j}}{\partial x_{i}}\right)\left(\frac{\partial u_{i}'}{\partial x_{j}}+\frac{\partial u_{j}'}{\partial x_{i}}\right)\dd^{3}\mathbf{x},
}\end{displaymath}%DIFAUXCMD
\DIFdelend \DIFaddbegin \begin{equation*}
    \DIFadd{\delta J=\int_{M}\delta\lambda\frac{\partial u_{i}}{\partial x_{i}}\frac{\partial u_{j}'}{\partial x_{j}}\dd^{3}\mathbf{x}+\frac{1}{2}\int_{M}\delta\mu\left(\frac{\partial u_{i}}{\partial x_{j}}+\frac{\partial u_{j}}{\partial x_{i}}\right)\left(\frac{\partial u_{i}'}{\partial x_{j}}+\frac{\partial u_{j}'}{\partial x_{i}}\right)\dd^{3}\mathbf{x},
}\end{equation*}\DIFaddend 
where we have set $\delta\sigma=0$ as that term has been dealt with. From this expression, we can now read off the desired sensitivity kernels
\DIFdelbegin \begin{eqnarray*}
    \DIFdel{K_{\lambda}=\frac{\partial u_{i}}{\partial x_{i}}\frac{\partial u_{j}'}{\partial x_{j}}, }\\
    \DIFdel{K_{\mu}=\left(\frac{\partial u_{i}}{\partial x_{j}}+\frac{\partial u_{j}}{\partial x_{i}}\right)\left(\frac{\partial u_{i}'}{\partial x_{j}}+\frac{\partial u_{j}'}{\partial x_{i}}\right).
}\end{eqnarray*}%DIFAUXCMD
\DIFdelend \DIFaddbegin \begin{gather*}
    \DIFadd{K_{\lambda}=\frac{\partial u_{i}}{\partial x_{i}}\frac{\partial u_{j}'}{\partial x_{j}}, }\\
    \DIFadd{K_{\mu}=\frac{1}{2}\left(\frac{\partial u_{i}}{\partial x_{j}}+\frac{\partial u_{j}}{\partial x_{i}}\right)\left(\frac{\partial u_{i}'}{\partial x_{j}}+\frac{\partial u_{j}'}{\partial x_{i}}\right).
}\end{gather*}\DIFaddend 
As is typical with the adjoint method, we see that these sensitivity kernels depend on both the solution of the forward and adjoint problems. These two problems are identical in form, differing only in their force terms, and hence the overall cost of the calculations is that of two forward calculations.



\item The part of the Lagrangian that depends explicitly on the motion is
\DIFdelbegin \begin{displaymath}
    \DIFdel{-\frac{1}{2}G\int_{M}\frac{\rho\rho'}{\{[\varphi_{i}-\varphi_{i}'][ \varphi_{i}-\varphi_{i}']\}^{\frac{1}{2}}}\dd^{3}\mathbf{x}',
}\end{displaymath}%DIFAUXCMD
\DIFdelend \DIFaddbegin \begin{equation*}
    \DIFadd{\frac{1}{2}G\int_{M}\frac{\rho\rho'}{\{[\varphi_{i}-\varphi_{i}'][ \varphi_{i}-\varphi_{i}']\}^{\frac{1}{2}}}\dd^{3}\mathbf{x}',
}\end{equation*}\DIFaddend 
where it is understood that \DIFdelbegin \DIFdel{un-primed }\DIFdelend \DIFaddbegin \DIFadd{unprimed }\DIFaddend functions depend on $\mathbf{x}$ and primed ones on the dummy integration variable $\mathbf{x}'$. Perturbing the motion to \DIFdelbegin \DIFdel{$\boldsymbol{\varphi}+\delta\boldsymbol{\varphi}$ }\DIFdelend \DIFaddbegin \DIFadd{$\bphi+\delta\bphi$ }\DIFaddend and expanding to first-order accuracy we find
\begin{equation*}
    \int_{M}\frac{\delta\mathcal{L}}{\delta\varphi_{i}}\delta\varphi_{i}\dd^{3}\mathbf{x}=-\frac{1}{2}G\int_{M}\int_{M}\frac{\rho\rho'(\varphi_{i}-\varphi_{i}')(\delta\varphi_{i}-\delta\varphi_{i}')}{\{[\varphi_{j}-\varphi_{j}'][ \varphi_{j}-\varphi_{j}']\}^{\frac{3}{2}}}\dd^{3}\mathbf{x}'\dd^{3}\mathbf{x}.
\end{equation*}
Splitting this expression into two parts we find
\begin{align*}
    \int_{M}\frac{\delta\mathcal{L}}{\delta\varphi_{i}}\delta\varphi_{i}\dd^{3}\mathbf{x} = &-\frac{1}{2}G\int_{M}\int_{M}\frac{\rho\rho'(\varphi_{i}-\varphi_{i}')\delta\varphi_{i}}{\{[\varphi_{j}-\varphi_{j}'][ \varphi_{j}-\varphi_{j}']\}^{\frac{3}{2}}}\dd^{3}\mathbf{x}'\dd^{3}\mathbf{x} \\
    &+\frac{1}{2}G\int_{M}\int_{M}\frac{\rho\rho'(\varphi_{i}-\varphi_{i}')\delta\varphi_{i}'}{\{[\varphi_{j}-\varphi_{j}'][ \varphi_{j}-\varphi_{j}']\}^{\frac{3}{2}}}\dd^{3}\mathbf{x}'\dd^{3}\mathbf{x},
\end{align*}
but by interchanging the order of integration and relabelling dummy variables $\mathbf{x} \leftrightarrow \mathbf{x}'$, we note that the two terms are actually equal, and hence
\begin{equation*}
    \int_{M}\frac{\delta\mathcal{L}}{\delta\varphi_{i}}\delta\varphi_{i}\dd^{3}\mathbf{x}=-G\int_{M}\delta\varphi_{i} \left( \int_{M}\frac{\rho\rho'(\varphi_{i}-\varphi_{i}')}{\{[\varphi_{j}-\varphi_{j}'][ \varphi_{j}-\varphi_{j}']\}^{\frac{3}{2}}}\dd^{3}\mathbf{x}' \right) \dd^{3}\mathbf{x},
\end{equation*}
which implies
\begin{equation*}
    \frac{\delta\mathcal{L}}{\delta\varphi_{i}}=-G\rho\int_{M}\frac{\rho'(\varphi_{i}-\varphi_{i}')}{\{[\varphi_{j}-\varphi_{j}'][ \varphi_{j}-\varphi_{j}']\}^{\frac{3}{2}}}\dd^{3}\mathbf{x}'=\rho\gamma_{i},
\end{equation*}
as required.

 \item Considering first the continuity of the potential across the boundary, we require
\begin{equation*}
    [\phi]_{-}^{+}=0
\end{equation*}
across \DIFdelbegin \DIFdel{$r=b+sh_{1}(\theta,\phi)+\cdot\cdot\cdot$}\DIFdelend \DIFaddbegin \DIFadd{$r=b+sh_{1}(\theta,\varphi)+\cdots$}\DIFaddend , where $[\cdot]_{-}^{+}$ denotes the jump in a quantity across the relevant boundary in the direction of the outward unit normal. Expanding this to first-order accuracy we obtain
\DIFdelbegin \begin{displaymath}
    [\DIFdel{\phi_{0}+s\frac{d\phi_{0}}{dr}h_{1}+s\phi_{1}+\cdot\cdot\cdot}]\DIFdel{_{-}^{+}=0,
}\end{displaymath}%DIFAUXCMD
\DIFdelend \DIFaddbegin \begin{equation*}
    \DIFadd{\left[\phi_{0}+s\frac{\ddns\phi_{0}}{\ddns r}h_{1}+s\phi_{1}+\cdots\right]_{-}^{+}=0,
}\end{equation*}\DIFaddend 
across the reference boundary $r=b$. Using the continuity of $\phi_{0}$ and its normal derivative across the reference boundary, this reduces to
\begin{equation*}
    [\phi_{1}]_{-}^{+}=0,
\end{equation*}
for $r=b$. To deal with the continuity of the normal derivative, we first write
\DIFdelbegin \begin{displaymath}
    \DIFdel{\hat{\mathbf{n}}=\hat{\mathbf{r}}+s\hat{\mathbf{n}}_{1}+\cdot\cdot\cdot,
}\end{displaymath}%DIFAUXCMD
\DIFdelend \DIFaddbegin \begin{equation*}
    \DIFadd{\hat{\mathbf{n}}=\hat{\mathbf{r}}+s\hat{\mathbf{n}}_{1}+\cdots,
}\end{equation*}\DIFaddend 
for the outward unit normal to the boundary. The term $\hat{\mathbf{n}}_{1}$ can be readily expressed in terms of $h_{1}$, but we will see that this is not required. Similar to above, the continuity of $\hat{\mathbf{n}}\cdot\nabla\phi$ across the aspherical boundary can be expanded to
\DIFdelbegin \begin{displaymath}
    [\DIFdel{\hat{\mathbf{r}}\cdot\nabla\phi_{0}+s(\hat{\mathbf{r}}\cdot\nabla)(\hat{\mathbf{r}}\cdot\nabla\phi_{0})h_{1}+s\hat{\mathbf{n}}_{1}\cdot\nabla\phi_{0}+s\hat{\mathbf{r}}\cdot\nabla\phi_{1}+\cdot\cdot\cdot}]\DIFdel{_{-}^{+}=0,
}\end{displaymath}%DIFAUXCMD
\DIFdelend \DIFaddbegin \begin{equation*}
    [\DIFadd{\hat{\mathbf{r}}\cdot\nabla\phi_{0}+s(\hat{\mathbf{r}}\cdot\nabla)(\hat{\mathbf{r}}\cdot\nabla\phi_{0})h_{1}+s\hat{\mathbf{n}}_{1}\cdot\nabla\phi_{0}+s\hat{\mathbf{r}}\cdot\nabla\phi_{1}+\cdots}]\DIFadd{_{-}^{+}=0,
}\end{equation*}\DIFaddend 
across $r=b$. The zeroth-order term vanishes immediately. We also note that $\hat{\mathbf{n}}_{1}\cdot\nabla\phi_{0}$ is zero as $\hat{\mathbf{n}}_{1}$ must be orthogonal to $\hat{\mathbf{r}}$ for the perturbed normal to be a unit vector, and this term is proportional to the radial derivative of $\phi_{0}$ which is continuous. The first-order condition therefore becomes
\begin{equation*}
    [(\hat{\mathbf{r}}\cdot\nabla)(\hat{\mathbf{r}}\cdot\nabla\phi_{0})h_{1}+\hat{\mathbf{r}}\cdot\nabla\phi_{1}]_{-}^{+}=0,
\end{equation*}
at $r=b$. To simplify the first term, we recall that $\phi_{0}$ satisfies the spherically symmetric Poisson equation
\DIFdelbegin \begin{displaymath}
    \DIFdel{\frac{1}{r^{2}}\frac{d}{dr}\left(r^{2}\frac{d\phi_{0}}{dr}\right)=4\pi G\rho_{0}.
}\end{displaymath}%DIFAUXCMD
\DIFdelend \DIFaddbegin \begin{equation*}
    \DIFadd{\frac{1}{r^{2}}\frac{\ddns}{\ddns r}\left(r^{2}\frac{\ddns\phi_{0}}{\ddns r}\right)=4\pi G\rho_{0}.
}\end{equation*}\DIFaddend 
Using this result, we see that
\begin{equation*}
    (\hat{\mathbf{r}}\cdot\nabla)(\hat{\mathbf{r}}\cdot\nabla\phi_{0}) = \frac{\partial^{2}\phi_{0}}{\partial r^{2}}=4\pi G\rho_{0}-\frac{2}{r}\frac{\partial\phi_{0}}{\partial r},
\end{equation*}
which can be substituted into the jump condition to obtain
\begin{equation*}
    \left[\frac{\partial\phi_{1}}{\partial r}+4\pi G\rho_{0}h_{1}\right]_{-}^{+}=0,
\end{equation*}
where again the continuity of the normal derivative of $\phi_{0}$ has been applied.




\item We are told
\begin{equation*}
    \mathbf{u}(\mathbf{x},t) = \sum_{km}\left\{\int_{0}^{t}\frac{\sin[\omega_{k}(t-t')]}{\omega_{k}}\int_{M}\overline{S}_{ij}(\mathbf{x}',t')\frac{\partial|km\rangle^{*}}{\partial x'_{j}}(\mathbf{x}')\dd^{3}\mathbf{x}'\dd t'\right\}|km\rangle(\mathbf{x}),
\end{equation*}
noting that arguments have now been included for clarity. In the case of a moment tensor point source, we see immediately that
\begin{equation*}
    \int_{M}\overline{S}_{ij}(\mathbf{x}',t')\frac{\partial|km\rangle^{*}}{\partial x'_{j}}(\mathbf{x}')\dd^{3}\mathbf{x}' = M_{ij}\frac{\partial|km\rangle^{*}}{\partial x_{j}}(\mathbf{x}_{s})H(t'-t_{s}),
\end{equation*}
while the necessary time integral is readily evaluated for $t\ge t_{s}$:
\begin{align*}
    &\int_{0}^{t}\frac{\sin[\omega_{k}(t-t')]}{\omega_{k}}H(t'-t_{s})\dd t' = \int_{t_s}^{t}\frac{\sin[\omega_{k}(t-t')]}{\omega_{k}}\dd t' \\ &= \left[\frac{\cos[\omega_k(t-t')]}{\omega_k^2}\right]_{t_s}^t = \frac{1-\cos[\omega_{k}(t-t_{s})]}{\omega_{k}^{2}},
\end{align*}
and hence the eigenfunction expansion reduces to
\begin{equation*}
    \mathbf{u}(\mathbf{x},t)=\sum_{km}\frac{1-\cos[\omega_{k}(t-t_{s})]}{\omega_{k}^{2}}M_{ij}\frac{\partial|km\rangle^{*}}{\partial x_{j}}(\mathbf{x}_{s})|km\rangle(\mathbf{x}),
\end{equation*}
for $t\ge t_{s}$ and zero before this. This final expression depends on ten source parameters, assuming that Earth structure is known. The dependence on the moment tensor is linear, but that on the source \DIFdelbegin \DIFdel{locations }\DIFdelend \DIFaddbegin \DIFadd{location }\DIFaddend and source time is non-linear. There are several possible methods for solving the inverse problem. The most obvious is to define some sort of least-squares misfit between observed and predicted seismograms, and then optimise to find the \DIFdelbegin \DIFdel{best fitting }\DIFdelend \DIFaddbegin \DIFadd{best-fitting }\DIFaddend source parameters. To do this gradient-based methods could be applied, and here the derivatives with respect to the source parameters could be calculated directly with relative ease. Note that in this problem there are likely to be more data than model parameters, and hence it is likely to be \DIFdelbegin \DIFdel{over-determined}\DIFdelend \DIFaddbegin \DIFadd{overdetermined}\DIFaddend . As a result, regularisation would not be needed. Because the problem is non-linear, it is not certain that the optimal solution obtained would be a global minimum. Due to the size of the inverse problem and the low cost of implementing mode summation (in a spherically symmetric earth \DIFaddbegin \DIFadd{model }\DIFaddend at least) the application of Bayesian methods via something like MCMC would be feasible, and this could address the possible issue of local minima, while also providing a quantification of uncertainty.

\item Within the eigenvalue problem we take the test function equal to $\mathbf{s}$ and so obtain
\DIFdelbegin \begin{displaymath}
    \DIFdel{-\omega^{2}+i\omega\langle \mathbf{s}|\mathbf{W}|\mathbf{s}\rangle+\langle \mathbf{s}|\mathbf{H}|\mathbf{s}\rangle=0,
}\end{displaymath}%DIFAUXCMD
\DIFdelend \DIFaddbegin \begin{equation*}
    \DIFadd{-\omega^{2}+\ii\omega\langle \mathbf{s}|W|\mathbf{s}\rangle+\langle \mathbf{s}|H|\mathbf{s}\rangle=0,
}\end{equation*}\DIFaddend 
where we have used the normalisation condition. This is a quadratic equation for $\omega$, and its solution yields
\DIFdelbegin \begin{displaymath}
    \DIFdel{\omega=\frac{1}{2}i\langle \mathbf{s}|\mathbf{W}|\mathbf{s}\rangle\pm\sqrt{\langle \mathbf{s}|\mathbf{H}|\mathbf{s}\rangle-\frac{1}{4}\langle \mathbf{s}|\mathbf{W}|\mathbf{s}\rangle^{2}},
}\end{displaymath}%DIFAUXCMD
\DIFdelend \DIFaddbegin \begin{equation*}
    \DIFadd{\omega=\frac{1}{2}\ii\langle \mathbf{s}|W|\mathbf{s}\rangle\pm\sqrt{\langle \mathbf{s}|H|\mathbf{s}\rangle-\frac{1}{4}\langle \mathbf{s}|W|\mathbf{s}\rangle^{2}},
}\end{equation*}\DIFaddend 
as required. Note that \DIFdelbegin \DIFdel{that }\DIFdelend the Coriolis operator is anti-Hermitian, and hence \DIFdelbegin \DIFdel{$\langle \mathbf{s}|\mathbf{W}|\mathbf{s}\rangle$ }\DIFdelend \DIFaddbegin \DIFadd{$\langle \mathbf{s}|W|\mathbf{s}\rangle$ }\DIFaddend is an imaginary number. For this reason the expression for $\omega$ is real so long as
\DIFdelbegin \begin{displaymath}
    \DIFdel{\langle \mathbf{s}|\mathbf{H}|\mathbf{s}\rangle-\frac{1}{4}\langle \mathbf{s}|\mathbf{W}|\mathbf{s}\rangle^{2}>0,
}\end{displaymath}%DIFAUXCMD
\DIFdelend \DIFaddbegin \begin{equation*}
    \DIFadd{\langle \mathbf{s}|H|\mathbf{s}\rangle-\frac{1}{4}\langle \mathbf{s}|W|\mathbf{s}\rangle^{2}>0,
}\end{equation*}\DIFaddend 
with the contribution from the Coriolis term always being non-negative. We see that in the rotating eigenvalue problem the Coriolis force always contributes to the stability of the earth model. As with the non-rotating case, the question of stability reduces to the properties of the potential energy form which here includes an additional contribution from centrifugal forces. This latter part can be written explicitly as
\DIFdelbegin \begin{displaymath}
    \DIFdel{\int_{M}\rho\epsilon_{ijk}\epsilon_{klm}s_{i}^{*}\Omega_{j}\Omega_{l}s_{m}\dd^{3}\mathbf{x}
}\end{displaymath}%DIFAUXCMD
\DIFdelend \DIFaddbegin \begin{equation*}
    \DIFadd{\int_{M}\rho\epsilon_{ijk}\epsilon_{klm}s_{i}^{*}\Omega_{j}\Omega_{l}s_{m}\dd^{3}\mathbf{x},
}\end{equation*}\DIFaddend 
and using the identity $\epsilon_{ijk}\epsilon_{klm}=\delta_{il}\delta_{jm}-\delta_{im}\delta_{jl}$ this becomes
\DIFdelbegin \begin{displaymath}
    \DIFdel{\int_{M}\rho\epsilon_{ijk}\epsilon_{klm}s_{i}^{*}\Omega_{j}\Omega_{l}s_{m}\dd^{3}\mathbf{x}=\int_{M}\rho[|(\mathbf{s}\cdot\boldsymbol{\Omega})|^2-||\mathbf{s}||^2||\boldsymbol{\Omega}||^2]\dd^{3}\mathbf{x}.
}\end{displaymath}%DIFAUXCMD
\DIFdelend \DIFaddbegin \begin{equation*}
    \DIFadd{\int_{M}\rho\epsilon_{ijk}\epsilon_{klm}s_{i}^{*}\Omega_{j}\Omega_{l}s_{m}\dd^{3}\mathbf{x}=\int_{M}\rho[|\mathbf{s}\cdot\bm{\Omega}|^{2}-\|\mathbf{s}\|^{2}\|\bm{\Omega}\|^{2}]\dd^{3}\mathbf{x}.
}\end{equation*}\DIFaddend 
Applying the \DIFdelbegin \DIFdel{Cauchy-Schwarz }\DIFdelend \DIFaddbegin \DIFadd{Cauchy--Schwarz }\DIFaddend inequality we note that this term is always non-positive, and hence centrifugal forces are destabilising.

\end{enumerate}

\end{document}
